\chapter{Introduction}

\section{Background and Motivation}

Parking is one of the last everyday transactions that still expects the driver to
interact with dedicated equipment. A typical urban facility installs a barrier, a
ticket dispenser or an RFID reader, a payment kiosk and, increasingly, a set of
cameras for number-plate recognition. Each item adds capital cost, a maintenance
visit and an independent attack surface, and all of it must be commissioned before
the first vehicle can enter. The hardware exists because a car has no identity that
a server can verify. The driver, however, does: a smartphone can authenticate a
person to a backend, display a price and authorise a payment.

Bluetooth Low Energy (BLE) completes the picture at the gate. BLE advertisements
are cheap to broadcast, readable by any modern phone, and can carry a facility
identifier without a network connection. Research on BLE-beacon parking has shown
that proximity-based discovery and automated payment are practical on commodity
phones \cite{mackey2020}, and surveys of IoT parking report the same direction:
sensing and connectivity move into low-cost devices while the business logic moves
into a backend \cite{alam2023,biyik2021}. What most of that work optimises is
\emph{occupancy sensing} --- finding a free bay --- while the barrier itself remains
a conventional, locally controlled device.

\project{} takes the opposite emphasis. It is not an occupancy system; it is a gate
controller. One ESP32 advertises its facility over BLE. The React Native
application in the driver's phone discovers the gate, starts a parking session on
the backend, and then writes the session identifier to the gate. The gate asks the
backend whether that session may pass, and moves a servo only when the answer is
\code{open}. There is no camera, no RFID tag, no ticket and no vehicle sensor. The
phone replaces the ticket, and the backend replaces the local decision.

The Indian toll ecosystem shows both the promise and the limit of vehicle-bound
tags: the National Electronic Toll Collection (NETC) programme issues a FASTag
RFID sticker per vehicle and settles tolls through a central system
\cite{npci2024}. That model works because a toll is a fixed charge at a fixed
point. Parking is neither: the charge depends on how long the vehicle stayed, and
the exit decision depends on whether it has been paid. \project{} therefore attaches
identity to the \emph{person} and time to the \emph{server}, which is what makes
per-second billing possible.

\section{Problem Statement}

Design and build a parking gate system in which:

\begin{enumerate}
  \item a driver discovers a nearby facility without scanning a QR code, typing a
        code or sharing a location;
  \item every entry and exit is authorised by a server, so that a forged client
        cannot manufacture permission to open a gate;
  \item billing is proportional to the actual time parked and is frozen the moment
        payment succeeds, with no clock controlled by the phone or the gate;
  \item the parking state machine cannot be driven into an invalid configuration
        by replaying requests, double-tapping a button or abandoning a session;
  \item the complete system --- backend, mobile application and gate firmware ---
        can be built, run and verified on a developer machine before any hardware
        is flashed.
\end{enumerate}

\section{Objectives}

\begin{enumerate}
  \item \textbf{Frictionless discovery.} Advertise the facility identity from the
        gate itself over BLE and resolve facility details from the backend.
  \item \textbf{Server-authoritative gate control.} Make the firmware a courier
        that acts only on a backend verdict and fails closed when the backend is
        unreachable.
  \item \textbf{Auditable per-second billing.} Compute every amount from the
        backend clock and the facility's pricing policy, and never accept a time
        or a price from the client.
  \item \textbf{A safe session lifecycle.} Enforce one active session per user and
        per vehicle, abort sessions that never reach a gate, and define a verdict
        for every illegal transition in advance.
  \item \textbf{A complete three-surface deliverable.} Deliver the Spring Boot
        backend, the Expo/React Native application and the ESP32 Arduino firmware
        from one repository with one configuration model.
  \item \textbf{Reproducible verification.} Provide one command that builds the
        images, boots the real application, runs the automated suite, enforces a
        coverage gate and publishes reports.
\end{enumerate}

\section{Scope and Limitations}

The version-one scope is deliberately narrow. Table~\ref{tab:scope} separates what
is implemented from what is excluded, and lists the limitations that follow from
the design decisions rather than from missing work.

\begin{table}[htbp]
  \centering
  \small
  \begin{tabularx}{\textwidth}{@{}lX@{}}
    \toprule
    \textbf{In scope} & \textbf{Out of scope (version one)} \\
    \midrule
    Registration, login and profile restore with JWT &
      Administrator user role or an operator web console \\
    Vehicle CRUD with number plate and class (MC/LMV/HMV) &
      Real payment-gateway integration, refunds and wallets \\
    BLE facility discovery from the gate's advertisement &
      GPS, QR-code or manual-code discovery \\
    Backend-authorised entry and exit with reason codes &
      Multi-gate facilities or more than one servo per site \\
    Per-second billing with frozen fees and receipts &
      Minimum fees, grace periods, daily caps, occupancy tracking \\
    Session history with status filter and paging &
      Heartbeats, over-the-air firmware updates, remote configuration \\
    Device provisioning and remote open/keep-closed commands &
      BLE pairing or encryption, iOS background BLE \\
    Dockerised build, test workflow and coverage gate &
      Production TLS termination and certificate management \\
    \bottomrule
  \end{tabularx}
  \caption{Scope of the delivered system.}
  \label{tab:scope}
\end{table}

Four limitations are accepted by design and are not defects:

\begin{itemize}
  \item \textbf{Plaintext BLE.} The link carries a session identifier, not a
        credential, and the gate re-validates every request with the backend, so
        the worst outcome of writing a foreign identifier is a closed gate.
  \item \textbf{Mock payment.} The provider sits behind a two-method interface and
        always succeeds; no money moves. Replacing it is an implementation of the
        interface, not a redesign.
  \item \textbf{No sensors.} The gate closes on a fixed eight-second timer. It
        cannot detect a tailgating vehicle, and it cannot detect a jammed arm except
        by a bounded settle window. A mechanical push-through override is required
        hardware.
  \item \textbf{Foreground BLE.} Discovery and gate commands need the application in
        the foreground and the phone within a few metres of the gate.
\end{itemize}
